\documentclass{amsart}
\usepackage{amsmath,amssymb,amsthm,hyperref}

\theoremstyle{definition}
\newtheorem{definition}{Definition}
\newtheorem{thesis}{Thesis}
\theoremstyle{plain}
\newtheorem{lemma}{Lemma}
\newtheorem{proposition}{Proposition}
\newtheorem{theorem}{Theorem}
\newtheorem{corollary}{Corollary}
\newtheorem{remark}{Remark}

\begin{document}

\title{Against Instrumentalism: The Irreducible Theoretical Function of
Scientific Theories}
\maketitle

\begin{abstract}
We settle the question posed: whether \emph{instrumentalism}---the view that
the meaning and value of a scientific theory are \emph{exhausted} by its
predictive utility as an instrument---is the correct primary interpretation of
scientific theories. We argue that it is mistaken, and we are careful to argue
this against the \emph{strongest} forms of instrumentalism rather than against
a stipulatively weak caricature. The argument fixes precise definitions; then
establishes two asymmetries between the use of statements as instruments
(applied science) and as conjectures (pure science): an \emph{asymmetry of
truth-bearing} and an \emph{asymmetry of risk}. We take explicit care to meet
(a) constructive empiricism, which appraises theories by \emph{empirical
adequacy} rather than convenience; (b) the \emph{Duhem--Quine} holism that
threatens the falsifiability asymmetry; (c) the \emph{pessimistic
meta-induction}, which threatens the dependence of instrumental success on
approximate truth; and (d) \emph{eliminative} instrumentalism, which grants the
truth-talk but calls it dispensable. The resulting verdict is scoped:
instrumentalism is correct as an account of the derivative \emph{applied} use,
but mistaken as the \emph{primary} interpretation, because theories possess a
descriptive, truth-apt function that is \emph{indispensable}---it does
explanatory work, especially the projection of theories into untested domains,
that no convenience-vocabulary can underwrite. Finally we determine what the
theoretical (non-practical) function of predictions consists in: not the
practical anticipation of events, but the production of \emph{potential
falsifiers} of the relevant test-system---the points at which a theory is
staked against reality and thereby treated as a description of the world.
\end{abstract}

\section{Formalization}

We first eliminate ambiguity by stating exactly what is being compared and
what counts as a successful resolution. Because the chief danger in a problem
of this kind is to win by definition, we will, after each definition, flag the
loophole through which a sophisticated instrumentalist might escape, and
discharge it later by argument rather than by stipulation.

\begin{definition}[Statements and their two uses]\label{def:uses}
Let a \emph{statement} be any declarative sentence of a scientific theory
(a universal law such as ``all planets move on ellipses,'' a formula such as
an equation of motion, or a singular prediction such as ``the eclipse begins
at $t_0$''). A statement may be \emph{used} in two ways.
\begin{enumerate}
\item[(A)] \textbf{Instrumental use} (applied science). The statement is used
as a \emph{tool} or rule of computation: it is fed inputs (initial conditions,
measured quantities) and yields outputs (predictions) \emph{for a given
practical purpose}. On this use the statement is appraised by its
\emph{performance relative to that purpose}: as efficient/inefficient,
applicable/inapplicable, convenient/inconvenient.
\item[(C)] \textbf{Conjectural use} (pure science). The statement is used as a
\emph{tentative description of, or conjecture about, reality}: an assertion put
forward as a candidate for being \emph{true} or \emph{false} of the world,
independently of any practical end.
\end{enumerate}
\end{definition}

\begin{definition}[Three grades of instrumentalism]\label{def:positions}
We distinguish three theses, in increasing order of philosophical strength
(harder to refute):
\begin{enumerate}
\item[(I$_w$)] \emph{Weak (pragmatic) instrumentalism}: a theory's value is its
convenience for a practical purpose; appraisal is purpose-indexed.
\item[(I$_e$)] \emph{Eliminative instrumentalism} (Def.\ here corresponds to
clause (ii) below): truth- and refutation-talk about theories is admitted as
\emph{psychologically natural but theoretically dispensable surplus}, to be
explained away in favour of predictive bookkeeping.
\item[(I$_a$)] \emph{Empirical-adequacy instrumentalism} (constructive
empiricism, after van Fraassen): a theory's epistemic value is its
\emph{empirical adequacy}---the truth of \emph{all} its observable
consequences---belief being withheld from its unobservable content. Here a
\emph{non}-relational truth predicate is admitted, but \emph{restricted to
observables}.
\end{enumerate}
The \emph{instrumentalist family} $\mathsf{INS}=\{\mathrm I_w,\mathrm
I_e,\mathrm I_a\}$ shares one commitment: the \emph{primary} appraisal of a
theory does not include its being a true (or approximately true)
\emph{description}, including of unobservable structure.
\end{definition}

\begin{definition}[The theoretical (realist) view]\label{def:realist}
The \emph{theoretical (realist) view} is the negation of the shared commitment:
scientific statements have a function of type (C)---a genuinely theoretical,
truth-apt, descriptive function reaching to unobservable structure---that is
\emph{irreducible} to, and \emph{not dispensable in favour of}, any function of
type (A).
\end{definition}

\begin{definition}[Reduction and dispensability]\label{def:reduce}
A function $F$ of statements \emph{reduces} to a function $G$ if every appraisal
proper to $F$ can be paraphrased without loss into the vocabulary proper to $G$.
$F$ is \emph{irreducible} to $G$ if some appraisal essential to $F$ cannot be so
paraphrased.

Separately, an appraisal $P$ is \emph{dispensable} relative to a vocabulary $V$
if there is a surrogate $P^{*}$ definable in $V$ such that replacing $P$ by
$P^{*}$ throughout the practice of science leaves \emph{every} legitimate
inference of that practice intact; $P$ is \emph{indispensable} if no such
surrogate exists---i.e.\ if some legitimate inference is licensed by $P$ but not
by any $V$-surrogate.
\end{definition}

\begin{remark}[The two burdens]\label{rem:burdens}
Refuting instrumentalism \emph{as a primary interpretation} requires two
things, corresponding to the two notions in
Definition~\ref{def:reduce}. \emph{(Irreducibility)} against $\mathrm I_w$ and
$\mathrm I_a$: the truth-apt descriptive appraisal is not paraphrasable into
convenience-talk (resp.\ into observable-adequacy-talk). \emph{(Indispensability)}
against $\mathrm I_e$: even granting the truth-talk, it cannot be eliminated
without losing legitimate scientific inferences. We discharge both. This is the
direct repair of the standard objection that the argument ``begs the question''
by smuggling its conclusion into Definition~\ref{def:uses}: nothing below is
proved from Definition~\ref{def:uses} alone; each load-bearing premise is
defended against a named opposing school.
\end{remark}

The question to be settled is therefore the truth value of the following.

\begin{thesis}[To be refuted]\label{th:inst}
Some member of $\mathsf{INS}$ gives the correct \emph{primary} interpretation of
scientific theories: the conjectural function \textup{(C)} either reduces to, or
is dispensable in favour of, the instrumental function \textup{(A)} (possibly
augmented, as in $\mathrm I_a$, by an observable-adequacy predicate).
\end{thesis}

\section{First asymmetry: truth-bearing, and the constructive-empiricist reply}

\subsection{Weak instrumentalism}

\begin{lemma}[Pure instruments are not true or false]\label{lem:instr-notrue}
Under weak instrumental use \textup{(I$_w$)}, the appraisals proper to a
statement are \emph{purpose-indexed} performance predicates
(applicable/inapplicable, efficient/inefficient, convenient/inconvenient
relative to a fixed purpose). The unrelativized predicate \emph{true} is not
among them.
\end{lemma}

\begin{proof}
The claim is \emph{conditional}: \textbf{if} appraisal is purpose-indexed (the
defining mark of $\mathrm I_w$), \textbf{then} truth, being non-relational,
is not among the admitted appraisals---for any predicate definable from
purpose-indexed predicates inherits a purpose-index, whereas truth does not
vary with purpose (``true for the engineer, false for the physicist'' is not a
coherent assignment of \emph{truth}, only of \emph{usefulness}). The content of
the lemma is thus a structural fact about $\mathrm I_w$'s own vocabulary, not a
hidden assertion that instrumentalism is false. It is corroborated by the
existence of \emph{deliberately false but maximally efficient} tools (the
flat-Earth approximation for surveying; point-mass molecules in elementary
kinetic theory): these are excellent \emph{for} their purpose and literally
false, so usefulness and truth come apart, confirming that the two appraisals
are distinct rather than that one is being defined away.
\end{proof}

\subsection{Empirical-adequacy instrumentalism: closing the false-dichotomy
loophole}

The previous lemma leaves open the strongest line: perhaps the relevant
contrast is not ``convenience vs.\ truth'' but ``\emph{observable} truth vs.\
\emph{unobservable} truth.'' Constructive empiricism ($\mathrm I_a$) grants a
non-relational truth predicate for observable consequences and asks why we need
more. We must therefore show that the descriptive function (C) is not captured
by \emph{empirical adequacy} (truth-of-all-observable-consequences). We do this
by argument, not stipulation.

\begin{lemma}[Empirical adequacy underdetermines, and so cannot replace,
descriptive content]\label{lem:adequacy}
Let $\mathrm{EA}(T)$ abbreviate ``every observable consequence of $T$ is
true.'' Then the appraisal ``$T$ is (approximately) a true description of the
world, including its unobservable structure'' is not paraphrasable into
appraisals built from $\mathrm{EA}$.
\end{lemma}

\begin{proof}
Suppose for contradiction that descriptive truth reduces to empirical adequacy.
Two facts block the reduction.

\emph{(i) Massive underdetermination of structure by observables.} For any
empirically adequate $T$ there are, in general, distinct theories $T'$ with
incompatible \emph{unobservable} commitments and the \emph{same} observable
consequences, hence $\mathrm{EA}(T)\leftrightarrow\mathrm{EA}(T')$. (Examples:
Lorentzian ether electrodynamics vs.\ special relativity at the observational
level; constructively equivalent reformulations of a field theory.) If
descriptive truth reduced to $\mathrm{EA}$, then $T$ and $T'$ would be
descriptively \emph{on a par}; yet they make \emph{contradictory} claims about
what the world is like. A predicate that assigns the same value to two
mutually inconsistent descriptions cannot \emph{be} the descriptive appraisal,
since the descriptive appraisal must, on pain of triviality, sometimes
separate inconsistent descriptions. Contradiction.

\emph{(ii) Adequacy is parasitic on description.} $\mathrm{EA}(T)$ is defined
as truth of the \emph{observable consequences of $T$}; these consequences are
generated by \emph{deduction from $T$'s full (incl.\ unobservable) content
together with bridge principles}. One cannot state which observable predictions
are even \emph{at issue} without first deploying $T$ as a description from which
they follow. Thus $\mathrm{EA}$ presupposes the descriptive use (C) in its very
definition; it is therefore a \emph{measure of} descriptive content's
observable shadow, not a \emph{replacement} for descriptive content. A
parasitic predicate cannot be that on which it is parasitic.

Either fact alone refutes the supposed reduction; hence descriptive truth is
not paraphrasable into empirical-adequacy talk.
\end{proof}

\begin{lemma}[Conjectures are essentially truth-apt]\label{lem:conj-true}
Under conjectural use \textup{(C)} the predicate \emph{true/false} (of the
world, unobservables included) is an essential appraisal: to use a statement as
a conjecture about reality \emph{just is} to put it forward as a candidate for
matching reality.
\end{lemma}

\begin{proof}
Immediate from Definition~\ref{def:uses}(C): a tentative description of reality
is a claim that reality is thus-and-so, hence constitutively the kind of thing
that can match or fail to match reality. Strip truth-aptness away and only a
rule remains---use (A), not (C).
\end{proof}

\begin{proposition}[First irreducibility, against $\mathrm I_w$ and $\mathrm
I_a$]\label{prop:irred1}
The conjectural function \textup{(C)} does not reduce to the instrumental
function \textup{(A)}, nor to (A) augmented by an empirical-adequacy predicate.
\end{proposition}

\begin{proof}
By Lemma~\ref{lem:conj-true} the truth-apt descriptive appraisal is essential
to (C). By Lemma~\ref{lem:instr-notrue} it is not expressible in $\mathrm
I_w$'s purpose-indexed vocabulary, and by Lemma~\ref{lem:adequacy} it is not
expressible in $\mathrm I_a$'s empirical-adequacy vocabulary. Hence an appraisal
essential to (C) cannot be paraphrased into the strongest instrumentalist
vocabularies; by Definition~\ref{def:reduce}, (C) is irreducible to them.
\end{proof}

\section{Second asymmetry: risk, holism, and the governance of localization}

The first asymmetry already blocks \emph{reduction}. Instrumentalism's
characteristic rejoinder is that truth is metaphysical surplus and that the
only respectable virtue is \emph{testability}. We meet this on its own ground
but must now respect the \emph{Duhem--Quine} thesis: a clash with observation
refutes only a whole \emph{test-system}, never an isolated law. We therefore do
\emph{not} claim naive single-counterexample falsification. Instead we relocate
the asymmetry into the \emph{normative governance} of how a clash is handled.

\begin{definition}[Test-system and localization]\label{def:testsystem}
A \emph{test-system} is a conjunction $S = T\wedge A\wedge c$, where $T$ is the
target theory, $A$ the auxiliary hypotheses (instruments, bridge principles,
background), and $c$ the initial conditions. A clash with an observation $o$
shows $S$ is false (some conjunct is false) but does not, by itself, say which.
A \emph{localization} is a choice of which conjunct to revise so as to restore
consistency with $o$.
\end{definition}

\begin{definition}[Risk-preserving vs.\ risk-dissolving
localization]\label{def:adhoc}
A localization is \emph{risk-preserving} if the revised conjunct is itself put
forward as a \emph{further conjecture about reality} that is independently
testable and forbids new observations (it has potential falsifiers of its own,
Def.~\ref{def:falsifier}). It is \emph{risk-dissolving} (\emph{ad hoc}) if its
\emph{sole} function is to immunize $S$ against the present clash---it forbids
nothing new, is not independently testable, and reduces the empirical content
of $S$.
\end{definition}

\begin{lemma}[Defeasible, holistic refutability of conjectural use]\label{lem:conj-refute}
Under conjectural use (C), the test-system $S=T\wedge A\wedge c$ is exposed to
\emph{defeasible} refutation: an observed clash establishes $\neg S$, and the
governing norm of (C) requires that the resulting localization be
\emph{risk-preserving}. Persisting clashes, under risk-preserving
localization, force eventual abandonment of $T$ as false.
\end{lemma}

\begin{proof}
By Lemma~\ref{lem:conj-true} each conjunct of $S$ is asserted of reality, so
$S$ is true only if reality is as $S$ says; a clash establishes $\neg S$ by
modus tollens. Holism shows this does not by itself fix \emph{which} conjunct
fails---granted. But the conjectural \emph{use} carries a norm: every conjunct
remains a conjecture \emph{at risk}. Hence the legitimate response is to replace
the suspected conjunct by another conjecture that is itself testable
(risk-preserving), not to neutralize the clash by a content-reducing manoeuvre.
Under repeated independent clashes that survive risk-preserving relocation, the
blame cannot be perpetually shifted without exhausting independently-testable
auxiliaries; at that point $T$ itself is convicted. This is refutation in the
only form holism permits: defeasible, system-level, and content-increasing in
its repairs.
\end{proof}

\begin{lemma}[Pure instrumental use lacks the governing norm of
risk]\label{lem:instr-norefute}
Under purely instrumental use ($\mathrm I_w$), a clash licenses, with equal
standing, \emph{any} localization that restores convenient performance,
including \emph{risk-dissolving} restriction of range; nothing in the
instrumental vocabulary marks a risk-dissolving localization as defective.
\end{lemma}

\begin{proof}
By Lemma~\ref{lem:instr-notrue}, $\mathrm I_w$ appraises only purpose-indexed
performance. A clash means ``inefficient/inapplicable \emph{here}.'' The
canonical repair---``restrict the tool's range; it remains efficient
\emph{there}''---is fully satisfactory by $\mathrm I_w$'s lights, and it
\emph{reduces} the domain over which the tool stakes a claim. Since $\mathrm
I_w$ has no predicate for ``forbids something new'' or ``stakes a fresh claim on
reality,'' it cannot distinguish a risk-preserving from a risk-dissolving
localization: both merely re-optimize convenience. Hence the norm of
Definition~\ref{def:adhoc} is \emph{inexpressible} in $\mathrm I_w$.
\end{proof}

\begin{proposition}[Second irreducibility, holism-aware]\label{prop:irred2}
The conjectural function \textup{(C)} does not reduce to the instrumental
function \textup{(A)} on the ground of risk-governance. The realist's constant
recourse to auxiliary revision (Lakatos's protective belt, Kuhnian normal
science) does \emph{not} collapse the asymmetry.
\end{proposition}

\begin{proof}
The appraisal essential to (C) is not ``immune to single counterexamples''
(which holism refutes for everyone) but ``\emph{governed by the norm of
risk-preserving localization}'' (Lemma~\ref{lem:conj-refute}). This appraisal
is unavailable to (A) (Lemma~\ref{lem:instr-norefute}), which cannot even
express the ad-hoc/legitimate distinction of Definition~\ref{def:adhoc}.
The objection that realists also revise auxiliaries is thus answered: the
relevant contrast is not \emph{whether} one localizes (both do) but \emph{under
what norm}. The conjectural user is bound to keep truth-commitment at risk by
demanding independently testable, content-increasing repairs; the purely
instrumental user is bound by nothing beyond convenience. An appraisal
inexpressible in (A) is, by Definition~\ref{def:reduce}, a witness to
irreducibility.
\end{proof}

\section{Primacy I: dependence without the no-miracles fallacy}

To lift the verdict from ``(C) is present and irreducible'' to ``(C) is
\emph{primary},'' the natural premise is that instrumental success
\emph{depends on} the underlying description's being approximately true. Stated
flatly (``reliable iff approximately true''), this is the no-miracles argument
and is refuted by Laudan's pessimistic meta-induction: caloric, phlogiston, the
optical ether, Ptolemaic astronomy and Newtonian gravitation were all reliable
\emph{within range} yet rest on now-false descriptions. We therefore do
\emph{not} assert the flat premise. We assert a \emph{selective} (structural)
version, defended below.

\begin{lemma}[Selective dependence]\label{lem:selective}
Whenever a theory $T$ is a reliable predictive instrument across a domain
$D$, there is a \emph{selectable} part of $T$'s descriptive content---the
mathematical-structural relations and the ``working posits'' actually deployed
in deriving the successful predictions---whose \emph{approximate truth} is
required for, and is retained by successors as, the explanation of that
reliability. The reliability does not require the approximate truth of $T$'s
\emph{idle} (non-working) posits.
\end{lemma}

\begin{proof}[Justification]
The pessimistic meta-induction targets \emph{whole theories} and their idle
ontological posits; it does not touch the \emph{structural} content. In each of
Laudan's cases the structural relations and working posits responsible for the
successful predictions are \emph{recovered as limiting cases} by the successor:
Fresnel's equations (structural relations among optical amplitudes) survive the
demise of the ether and reappear in Maxwell's theory; Newtonian gravitation's
inverse-square structure survives as the weak-field, low-velocity limit of
general relativity; Carnot's structural relations survive the demise of
caloric. The general pattern (selective/structural realism, Worrall, Psillos):
\emph{structure and working posits are conserved across theory change, idle
posits are discarded.} Hence the correct dependence claim is not ``reliability
$\Rightarrow$ whole theory approximately true'' but ``reliability across $D$
$\Rightarrow$ the deployed structural/working content is approximately true of
$D$,'' which Laudan's list confirms rather than refutes (the survivors are
exactly the working structure).
\end{proof}

\begin{proposition}[Asymmetric dependence]\label{prop:depend}
The instrumental use presupposes the conjectural use, but not conversely.
\end{proposition}

\begin{proof}
By Lemma~\ref{lem:selective}, to trust $T$'s outputs as predictions over a
\emph{novel} stretch of $D$ is to rely on the approximate truth of $T$'s
deployed structural content over that stretch---a type-(C) commitment. Without
some such believed-true structural description, extrapolation of the tool beyond
the data already in hand would be unlicensed (see
Theorem~\ref{thm:indispensable}). Conversely a conjecture can be entertained,
tested, and refuted with no practical application in view
(Lemma~\ref{lem:conj-refute}; theoretical predictions are tested for their
truth long before, or without ever, being used). Thus (A) depends on (C) but not
(C) on (A): (C) is explanatorily and historically prior.
\end{proof}

\section{Primacy II: indispensability against eliminative instrumentalism}

The eliminative instrumentalist ($\mathrm I_e$) concedes the truth-talk but
calls it dispensable surplus: replace ``$T$ is approximately true'' by a
surrogate ``$T$ has been convenient/empirically adequate so far,'' and---he
claims---science proceeds unchanged. By Definition~\ref{def:reduce} we must
exhibit a \emph{legitimate inference} licensed by the truth-appraisal but by no
$\mathrm I_e$-surrogate. The decisive case is \emph{projection into an untested
domain}.

\begin{lemma}[Past-convenience surrogates do not project]\label{lem:noproject}
Let $P$ be the appraisal ``$T$'s deployed structural content is approximately
true (of the relevant kind of system).'' Let $P^{*}$ be any surrogate definable
in $\mathrm I_e$'s vocabulary from $T$'s track record (e.g.\ ``$T$ was
convenient/empirically adequate on the tested domain $D_0$''). Then $P^{*}$ does
not license, while $P$ does license, the inference to $T$'s reliability on a
\emph{new} domain $D_1\not\subseteq D_0$ of the same kind.
\end{lemma}

\begin{proof}
$P^{*}$ is, by construction, a summary of performance on $D_0$. Any inference
from ``performed well on $D_0$'' to ``will perform well on $D_1$'' is a bare
enumerative induction over \emph{instances of use}; it is not underwritten by
$P^{*}$ itself, which is silent about $D_1$. The only thing that bridges $D_0$
and $D_1$ is a claim that the \emph{same structure obtains} in both---i.e.\ a
type-(C) descriptive claim that $T$ \emph{truly describes} the relevant kind of
system, of which $D_0$ and $D_1$ are instances. This is exactly $P$. Concretely:
an engineer extends a stress formula to a novel alloy geometry, or a physicist
applies a field equation to a regime never measured, \emph{because the theory is
believed to describe the underlying mechanism}, not merely because it was handy
before. Strip the descriptive belief and the projection loses its warrant: past
convenience entails nothing about an unsampled domain. Hence $P$ licenses an
inference $P^{*}$ cannot.
\end{proof}

\begin{theorem}[Indispensability of the descriptive appraisal]\label{thm:indispensable}
The truth-apt descriptive appraisal (C) is \emph{indispensable} in the sense of
Definition~\ref{def:reduce}: replacing it by any instrumentalist surrogate
forfeits a legitimate and ubiquitous scientific inference---the projection of a
theory into untested domains.
\end{theorem}

\begin{proof}
Projection into untested domains is a legitimate, indeed routine, inference of
both pure and applied science (every new application of an old theory is one).
By Lemma~\ref{lem:noproject} this inference is licensed by $P$ but by no
$\mathrm I_e$-surrogate. By Definition~\ref{def:reduce} the appraisal is
therefore indispensable: it cannot be eliminated salva the inferential
practice. This refutes the eliminative option $\mathrm I_e$, which claimed
exactly such eliminability.
\end{proof}

\section{The main theorem}

\begin{theorem}[Instrumentalism is mistaken as a primary
interpretation]\label{thm:main}
Thesis~\textup{\ref{th:inst}} is false against every member of $\mathsf{INS}$.
Scientific theories possess a theoretical (descriptive, truth-apt, structurally
committed, risk-governed) function that is both \emph{irreducible} to and
\emph{indispensable} relative to their instrumental function, and on which the
instrumental function depends. Therefore no member of $\mathsf{INS}$ gives the
correct primary interpretation of scientific theories.
\end{theorem}

\begin{proof}
\emph{Irreducibility.} Against $\mathrm I_w$ and $\mathrm I_a$,
Proposition~\ref{prop:irred1} (truth-bearing) and
Proposition~\ref{prop:irred2} (risk-governance) show (C) is not paraphrasable
into their vocabularies.

\emph{Indispensability.} Against $\mathrm I_e$,
Theorem~\ref{thm:indispensable} shows the descriptive appraisal cannot be
eliminated without losing the projection inference.

\emph{Primacy.} By Proposition~\ref{prop:depend} the instrumental use depends on
the conjectural use and not conversely; a function on which the other depends,
and which is both irreducible and indispensable, cannot be the \emph{secondary}
one. Moreover an interpretation is correct only if it preserves the appraisals
science actually makes: science distinguishes \emph{true from false structural
description} (Prop.~\ref{prop:irred1}), \emph{risk-preserving from ad hoc
repair} (Prop.~\ref{prop:irred2}), and \emph{warranted from unwarranted
projection} (Thm.~\ref{thm:indispensable}); the purely instrumental
vocabularies collapse all three. Hence instrumentalism, while a correct account
of the derivative \emph{applied} use (A), is mistaken as the \emph{primary}
interpretation of scientific theories.
\end{proof}

\begin{remark}[Scope of the verdict]\label{rem:scope}
The verdict refutes the \emph{primacy} claim shared by $\mathrm I_w,\mathrm
I_e,\mathrm I_a$: none can be the correct primary interpretation, because each
either cannot express, or cannot dispense with, the descriptive function. It is
compatible with the modest constructive-empiricist thesis that \emph{belief} in
unobservables is epistemically optional; what is refuted is the stronger claim
that the descriptive, truth-apt function is not the \emph{primary} role of a
theory. Where the realist and the constructive empiricist genuinely part
company---over the \emph{rational obligatoriness} of belief in unobservable
structure---lies beyond the present, conceptual, question; the present question
concerns the \emph{function} of theories, and on that question the descriptive
function is shown primary.
\end{remark}

\section{The theoretical function of predictions}

The problem asks what the theoretical (non-practical) function of
\emph{predictions} consists in on the view just established.

\begin{definition}[Potential falsifier]\label{def:falsifier}
A \emph{potential falsifier} of a test-system $S=T\wedge A\wedge c$ is a
singular observation-statement $f$ logically incompatible with $S$, so that,
were $f$ observed, $S$ would be refuted in the sense of
Lemma~\ref{lem:conj-refute}. The class of potential falsifiers is the class of
observations $S$ \emph{forbids}.
\end{definition}

\begin{proposition}[Predictions as potential
falsifiers]\label{prop:pred-falsifier}
A prediction derived from $T$ together with auxiliaries $A$ and initial
conditions $c$ is, on the theoretical view, an instruction for generating a
potential falsifier of the test-system $S=T\wedge A\wedge c$: it specifies in
advance an observable outcome whose negation, if observed, would refute $S$
(and, under risk-preserving localization, eventually $T$). Its theoretical
(non-practical) function is to \emph{expose $S$ to refutation}---to convert
$T$'s descriptive claim into a definite confrontation with reality.
\end{proposition}

\begin{proof}
Let $T$ be used conjecturally and deduce from $S=T\wedge A\wedge c$ a singular
prediction $p$: ``the outcome will be $O$.'' Then $\bar p$: ``the outcome is not
$O$'' is logically incompatible with $S$; observing $\bar p$ refutes $S$ by
Lemma~\ref{lem:conj-refute}. Hence $\bar p$ is a potential falsifier of $S$
(Def.~\ref{def:falsifier}), and issuing $p$ is the act of singling out, in
advance, an observation $S$ forbids. The point of issuing $p$ in pure science is
not to \emph{use} the outcome for an external purpose but to \emph{stake $S$'s
descriptive claim on it}. This is a function of type (C): measured not by
convenience but by what it would teach us about the truth of the structural
content of $T$ (the conjunct kept at risk by the governing norm,
Def.~\ref{def:adhoc}).
\end{proof}

\begin{corollary}[The non-practical point of prediction]\label{cor:why-predict}
On the resulting view the theoretical function of a prediction is twofold and
purely epistemic:
\begin{enumerate}
\item \textbf{As a test.} A prediction is the cutting edge by which a theory is
brought into contact with reality; the more precise and the more
\emph{improbable in advance} the outcome, the larger the class of potential
falsifiers it activates, hence the more severely $T$ is tested.
\item \textbf{As a bearer of empirical content.} Because a refutable prediction
\emph{forbids} something, it has content (it excludes possible states of the
world). A theory's descriptive, explanatory grasp is measured by the totality of
what it forbids; predictions are the concrete tokens of that content.
\end{enumerate}
In neither role is the prediction valued for practical use; in both it serves
the theoretical end of learning whether, and how far, the theory's description
of reality is true.
\end{corollary}

\begin{proof}
Both clauses follow from Proposition~\ref{prop:pred-falsifier} with
Definition~\ref{def:falsifier}. For (1): the falsifier class activated by $p$ is
the set of outcomes incompatible with $p$ given $A,c$; the less antecedently
probable $O$, the larger that class, hence the greater the risk---the measure of
test severity. For (2): a statement has empirical content to the extent that it
has potential falsifiers; by Proposition~\ref{prop:pred-falsifier} a prediction
has a nonempty such class, so it carries content, and the aggregate of forbidden
observations is the theory's empirical content---the precise sense in which it
says something about the world. Neither measure refers to a practical purpose;
both are type-(C).
\end{proof}

\section{Conclusion}

\begin{theorem}[Resolution]\label{thm:resolution}
Instrumentalism is correct only as an account of the derivative, \emph{applied}
use of scientific statements as tools, and is \emph{mistaken} as the
\emph{primary} interpretation of scientific theories---this against weak,
eliminative, and empirical-adequacy instrumentalism alike. Theories possess an
irreducible and indispensable theoretical (descriptive, truth-apt, structurally
committed, risk-governed) function, on which their instrumental use itself
depends. Correspondingly, the theoretical (non-practical) function of
predictions is not the anticipation of events for use, but the generation of
potential falsifiers---the points at which the test-system, and through it the
theory, is staked against reality and thereby treated, and tested, as a
description of the world.
\end{theorem}

\begin{proof}
The first two sentences are Theorem~\ref{thm:main} together with
Remark~\ref{rem:scope}; the irreducibility, indispensability and dependence
clauses are Propositions~\ref{prop:irred1}, \ref{prop:irred2},
\ref{prop:depend} and Theorem~\ref{thm:indispensable}. The last sentence is
Proposition~\ref{prop:pred-falsifier} with Corollary~\ref{cor:why-predict}.
Combining them yields the stated resolution.
\end{proof}

\end{document}
